%% SSEE-V3.6 — Apéndice Técnico: Tensor EMI y Ecuaciones de Friedmann
%% Autor: Mike Edison Almeida Vallejo
%% Derivación directa desde T_{\mu\nu} (complementa SSEE_EFT_section.tex)
%% Correcciones aplicadas: (1) nota física Ec.7, (2) "traza espacial" en lugar de "ii",
%%   (3) ecuación intermedia explicitada, (4) conclusión fenomenológica completada.

\section*{Apéndice Técnico: Derivación del Tensor de Energía-Impulso SSEE
          y Ecuaciones de Friedmann Modificadas}

En este apéndice se formaliza la introducción de la constante de retención
geométrica ($\mathrm{KAL}_{0}$) en la dinámica de expansión cósmica.
Dado que la viscosidad \textit{bulk} introduce irreversibilidad temporal
---los procesos disipativos aumentan la entropía--- no pueden derivarse de
una acción variacional clásica, que es invariante bajo inversión temporal.
Por eso, en lugar de forzar $\Pi$ en una Lagrangiana reversible, postulamos
$T_{\mu\nu}^{\mathrm{SSEE}}$ directamente, adoptando una \emph{variante}
del formalismo de Eckart para fluidos no perfectos en la que el coeficiente
de viscosidad $\zeta$ es proporcional a la tasa de expansión $H$ (extensión
respecto al formalismo de Eckart estándar, donde $\zeta$ es función de
$\rho$ o constante \citep{Weinberg1971}).

% ─────────────────────────────────────────────────────────────────────────────
\subsection*{1. El Tensor de Energía-Impulso Efectivo SSEE}

El universo se modela como un fluido homogéneo e isotrópico.
La contribución del sector oscuro geométrico SSEE se describe mediante un
tensor que incluye una presión termodinámica estándar y una presión viscosa
(\textit{bulk viscosity}):
%
\begin{equation}
  T_{\mu\nu}^{\mathrm{SSEE}}
    = \bigl(\rho_{\mathrm{SSEE}} + p_{\mathrm{eff}}\bigr)u_\mu u_\nu
      + p_{\mathrm{eff}}\,g_{\mu\nu},
  \label{eq:app_Tmunu}
\end{equation}
%
donde la presión efectiva es $p_{\mathrm{eff}} = p_{\mathrm{SSEE}} + \Pi$.
El sector de energía oscura obedece una ecuación de estado dictada
algebraicamente por la geometría, $p_{\mathrm{SSEE}} = w_{0}\rho_{\mathrm{SSEE}}$
(con $w_{0} \approx -0.840$).

La presión viscosa $\Pi$ se define, siguiendo la estructura de Eckart,
como proporcional a la expansión volumétrica del fluido
($\nabla_{\!\mu}u^{\mu} = 3H$ en FLRW):
%
\begin{equation}
  \Pi = -\zeta\,\nabla_{\!\mu}u^{\mu} = -3\zeta H.
  \label{eq:app_Pi_def}
\end{equation}

% ─────────────────────────────────────────────────────────────────────────────
\subsection*{2. La Condición de Retención SSEE (Sin Parámetros Fenomenológicos Adicionales)}

A diferencia de los modelos cosmológicos estándar que tratan a $\zeta$ como
un parámetro libre a ajustar con datos, el marco SSEE postula que todo el
sector viscoso está codificado en la constante adimensional
$\mathrm{KAL}_{0} = \beta + \pi \approx 5.521$, derivada algebraicamente
de $\varphi$ y $\pi$ sin ningún ajuste numérico.
Para satisfacer la consistencia dimensional exacta
($[\zeta] = E^{3}$ en unidades naturales), acoplamos $\mathrm{KAL}_{0}$
a la tasa de expansión y a la constante de gravitación:
%
\begin{equation}
  \zeta = \mathrm{KAL}_{0}\,\frac{H}{8\pi G}.
  \label{eq:app_zeta}
\end{equation}
%
Todo factor en~(\ref{eq:app_zeta}) es o bien una constante fundamental ($G$)
o bien un invariante algebraico SSEE ($\mathrm{KAL}_{0}$, $H$ determinado
por las ecuaciones de fondo).
En particular, no se introduce ningún parámetro fenomenológico adicional
$\kappa$ del tipo $\zeta = \kappa\,\mathrm{KAL}_{0}\,H/8\pi G$;
el único parámetro efectivo del sector disipativo es $\mathrm{KAL}_{0}$,
cuyo valor está fijado algebraicamente antes de comparar con observaciones.

Sustituyendo~(\ref{eq:app_zeta}) en~(\ref{eq:app_Pi_def}):
%
\begin{equation}
  \Pi
    = -3\!\left(\mathrm{KAL}_{0}\,\frac{H}{8\pi G}\right)\!H
    = -\frac{3\,\mathrm{KAL}_{0}\,H^{2}}{8\pi G}.
  \label{eq:app_Pi}
\end{equation}
%
Nótese que $\Pi < 0$ (presión viscosa negativa), consistente con la
retención disipativa de energía a alto $H$ \citep{Weinberg1971, Brevik2005}.

% ─────────────────────────────────────────────────────────────────────────────
\subsection*{3. Derivación de las Ecuaciones de Friedmann}

Al inyectar $T_{\mu\nu}^{\mathrm{SSEE}}$ en las ecuaciones de campo de
Einstein ($G_{\mu\nu} = 8\pi G\,T_{\mu\nu}$) sobre la métrica FLRW
$ds^{2} = -dt^{2} + a^{2}(t)\,d\mathbf{x}^{2}$, se obtienen tres leyes
de expansión.

\paragraph{I. Ecuación de la Energía (componente $00$):}
La densidad total determina la tasa de expansión:
%
\begin{equation}
  H^{2} = \frac{8\pi G}{3}\bigl[\rho_{m} + \rho_{r} + \rho_{\mathrm{SSEE}}\bigr].
  \label{eq:app_Friedmann_H}
\end{equation}

\paragraph{II. Ecuación de la Aceleración (traza de las componentes espaciales):}
A partir de la ecuación de Raychaudhuri,
$\ddot{a}/a = -(4\pi G/3)(\rho + 3p_{\mathrm{eff}})$, se obtiene:
%
\begin{equation}
  \frac{\ddot{a}}{a}
    = -\frac{4\pi G}{3}
      \bigl[\rho_{m} + 2\rho_{r}
            + \rho_{\mathrm{SSEE}}(1 + 3w_{0}) + 3\Pi\bigr].
  \label{eq:app_acc_presubst}
\end{equation}
%
Sustituyendo $3\Pi = -9\,\mathrm{KAL}_{0}H^{2}/8\pi G$, llegamos a la
ecuación maestra de aceleración SSEE:
%
\begin{equation}
  \frac{\ddot{a}}{a}
    = -\frac{4\pi G}{3}
      \!\left[\rho_{m} + 2\rho_{r}
        + \rho_{\mathrm{SSEE}}(1 + 3w_{0})
        - \frac{9\,\mathrm{KAL}_{0}\,H^{2}}{8\pi G}
      \right].
  \label{eq:app_Friedmann_a}
\end{equation}

\textit{Nota física:}
El término viscoso $\Pi < 0$ actúa como presión efectiva negativa adicional.
Su contribución al lado derecho de~(\ref{eq:app_Friedmann_a}) es
$-(4\pi G/3)\cdot(-9\,\mathrm{KAL}_{0}H^{2}/8\pi G)
 = +3\,\mathrm{KAL}_{0}H^{2}/2 > 0$,
es decir, una contribución \emph{positiva} a $\ddot{a}/a$ que modera la
deceleración gravitacional sin requerir partículas exóticas ni ajuste fino.

\paragraph{Derivación del coeficiente 9.}
De~(\ref{eq:app_Pi}): $\Pi = -3\,\mathrm{KAL}_{0}H^{2}/8\pi G$,
por lo tanto $3\Pi = -9\,\mathrm{KAL}_{0}H^{2}/8\pi G$.
El factor 9 proviene de $3\times|\Pi|/H^{2}
= 3\times(3\,\mathrm{KAL}_{0}/8\pi G)$;
un factor 3 en lugar de 9 correspondería a incluir $\Pi$ una sola vez
en vez de como la traza del tensor de estrés viscoso.

% ─────────────────────────────────────────────────────────────────────────────
\subsection*{4. Ecuación de Continuidad y Generación Dinámica}

La conservación del tensor de energía-impulso
($\nabla_{\!\mu}T^{\mu\nu} = 0$) para el sector SSEE produce:
%
\begin{equation}
  \dot{\rho}_{\mathrm{SSEE}}
    + 3H\bigl(\rho_{\mathrm{SSEE}} + p_{\mathrm{eff}}\bigr) = 0.
  \label{eq:app_cons_raw}
\end{equation}
%
Separando $p_{\mathrm{eff}} = w_{0}\rho_{\mathrm{SSEE}} + \Pi$:
%
\begin{equation}
  \dot{\rho}_{\mathrm{SSEE}}
    + 3H\bigl[\rho_{\mathrm{SSEE}}(1 + w_{0}) + \Pi\bigr] = 0.
  \label{eq:app_cons_split}
\end{equation}
%
Sustituyendo la definición SSEE de $\Pi$ (Ec.~\ref{eq:app_Pi}),
obtenemos la ecuación de continuidad modificada:
%
\begin{equation}
  \dot{\rho}_{\mathrm{SSEE}}
    + 3H\,\rho_{\mathrm{SSEE}}\,(1 + w_{0})
    = \frac{9\,\mathrm{KAL}_{0}\,H^{3}}{8\pi G}.
  \label{eq:app_continuity}
\end{equation}

\textit{Conclusión fenomenológica:}
En el universo temprano ($H \gg H_{0}$), el término disipativo
$9\,\mathrm{KAL}_{0}H^{3}/8\pi G$ domina el lado derecho de
(\ref{eq:app_continuity}), actuando como una fuente que inyecta energía
en el fluido oscuro y modifica su ecuación de estado efectiva.
Este mecanismo de \emph{retención geométrica} es la base física del
régimen CMB del modelo SSEE: en el límite $H \to \infty$ el término
$\propto H^{3}$ actúa como una fuente inercial que reemplaza parcialmente
la contribución de la materia oscura fría, permitiendo
$\Omega_{m,\mathrm{eff}} = 1 - T_{r}/M_{v} \approx 0.160$
compatible con la escala acústica del CMB a través del factor MIRA
$\mathcal{M} = 2 - 1/\varphi^{2} \approx 1.999$ (Papel~3, §2.3).
Una derivación cuantitativa de $\Omega_{m,\mathrm{eff}}(z)$ y de la
función de crecimiento de estructuras, que conecta este límite con los
datos Planck~PR4, se presenta en Papel~3.
En el universo tardío ($H \approx H_{0}$) el término cúbico es
despreciable y el sector SSEE se comporta esencialmente como un fluido
k-essence con $w_{0} \approx -0.840$, en acuerdo con los datos DESI~DR2
a $0.05\sigma$ (Papel~2).
